%==============================================================================
\section{Neural Networks --- Mạng nơ-ron và học máy}
%==============================================================================

\begin{ynghia}[title={Nguồn}]
\tsurl{https://www.tutorialspoint.com/machine_learning/machine_learning_artificial_neural_networks.htm}.
\end{ynghia}

\begin{dongco}[title={Hai khái niệm thường đi cùng nhau}]
Học máy và mạng nơ-ron là hai công nghệ quan trọng trong AI.
Chúng thường được dùng cùng nhau nhưng \textbf{không} trùng nhau.
Bài này làm rõ khái niệm từng bên và so sánh sự khác biệt.
\end{dongco}

\subsection{Học máy là gì?}

\begin{dinhnghia}[title={Machine Learning}]
Machine Learning là lĩnh vực khoa học máy tính giúp hệ máy tính ``hiểu'' dữ liệu theo cách gần với con người.
Nói đơn giản: ML là một dạng AI trích xuất pattern từ dữ liệu thô bằng thuật toán hoặc phương pháp.
\end{dinhnghia}

\begin{phuongphap}[title={Ba nhóm học máy (theo mức giám sát)}]
\begin{itemize}
  \item \textbf{Supervised learning}: huấn luyện trên dữ liệu có nhãn cho phân loại và hồi quy.
  Ví dụ thuật toán: linear regression, K-nearest neighbors, decision trees, random forest, \ldots
  \item \textbf{Unsupervised learning}: huấn luyện trên dữ liệu không nhãn; dùng cho clustering, association rule mining, giảm chiều.
  Ví dụ: K-means, thuật toán Apriori, \ldots
  \item \textbf{Reinforcement learning}: agent tương tác môi trường qua hành động và quan sát kết quả; học dựa trên thưởng/phạt.
  Ví dụ: Q-learning, policy gradient, Monte Carlo, \ldots
\end{itemize}
\end{phuongphap}

\begin{luuy}[title={Đối chiếu}]
Ba kiểu học: Mục~\ref{sec:bon-kieu-hoc}; bài này nhấn góc nhìn \textbf{mạng nơ-ron}.
\end{luuy}

\subsection{Mạng nơ-ron là gì?}

\begin{dinhnghia}[title={Neural Networks}]
Mạng nơ-ron là \textbf{một loại} thuật toán học máy, lấy cảm hứng từ cấu trúc não người.
Chúng mô phỏng cách não hoạt động bằng các lớp node (nơ-ron nhân tạo) liên kết với nhau.
Mỗi nơ-ron nhận đầu vào từ lớp trước và tạo đầu ra; quá trình lặp qua các lớp cho đến đầu ra cuối.
\end{dinhnghia}

\begin{ynghia}[title={Ứng dụng}]
Mạng nơ-ron dùng cho nhiều tác vụ: nhận dạng ảnh, giọng nói, NLP, dự báo.
Đặc biệt phù hợp khi xử lý dữ liệu phức tạp hoặc nhận diện pattern trong dữ liệu.
\end{ynghia}

\begin{vidu}[title={Mô hình ANN tổng quát (theo nguồn)}]
\tsimg{09_ai_neural_deep/artificial_neural_network_model.jpg}
\begin{center}\footnotesize Nguồn: \tsurl{https://www.tutorialspoint.com/machine_learning/images/artificial_neural_network_model.jpg}\end{center}
Sơ đồ trên mô tả mô hình mạng nơ-ron nhân tạo tổng quát và quá trình xử lý qua các lớp.
\end{vidu}

\subsection{Machine Learning và Neural Networks}

\begin{ynghia}[title={Bốn điểm so sánh chính (theo nguồn)}]
\begin{enumerate}
  \item ML là \textbf{nhóm rộng} gồm nhiều thuật toán, trong đó có mạng nơ-ron.
  \item ML dùng cho nhiều tác vụ; mạng nơ-ron đặc biệt mạnh với dữ liệu phức tạp và pattern tinh vi mà thuật toán khác có thể bỏ sót.
  \item Mạng nơ-ron thường cần \textbf{nhiều dữ liệu} và sức mạnh tính toán (ví dụ GPU) hơn nhiều thuật toán ML cổ điển.
  \item Mạng nơ-ron có thể cho dự đoán/quyết định rất chính xác nhưng đôi khi \textbf{khó giải thích} hơn các thuật toán ML khác (hộp đen).
\end{enumerate}
\end{ynghia}

\begin{luuy}[title={Giải thích mô hình}]
Cách mạng nơ-ron ra quyết định không luôn minh bạch, nên khó truy vết lý do đến kết luận --- đây là thách thức quan trọng trong triển khai thực tế (liên quan interpretability ở các bài trước).
\end{luuy}
